\section{Introduction}\label{sec:introduction}
% statements-of-record rows resolved here: none

The Birch--Swinnerton-Dyer formula~\citep{BirchSwinnertonDyer1965,Tate1966bsd} is usually first met as a scalar
identity.  For an elliptic curve $E/\Q$ of rank $r$, the standard leading
coefficient form compares
\[
  \frac{L^{(r)}(E,1)}{r!}
  \quad\text{with}\quad
  \Omega_E\,\Reg(E/\Q)\,
  \frac{\#\Sha(E/\Q)\,\Tam(E)}{\#E(\Q)_{\tors}^{\,2}} .
\]
This product is already visibly composite: an analytic leading term and
period, a height regulator, a finite Tamagawa--torsion--$\Sha$ factor, local
$p$-adic data, and the determinant-line conventions that make all of these
terms live in a common orientation.  The question addressed in this paper is
not whether the scalar BSD identity is true.  It is the more local diagnostic
question: when one separates the Bloch--Kato fundamental line~\citep{BlochKato1990,FontainePerrinRiou1994} into these
typed pieces, which pieces carry enough native data to recover their intended
readout, and where is information lost by passing to a public scalar shadow?

The word ``typed'' is used here in its elementary sense.  We keep the five
classes of data in five labelled slots rather than immediately multiplying
them into one number.  The height/regulator slot contains the Mordell--Weil
lattice and the Neron--Tate pairing; the analytic/period slot contains the
leading Taylor coefficient and the real period; the finite slot contains
Tamagawa factors, torsion, the $p$-primary $\Sha$ data, and Cassels--Tate
pairings; the $p$-adic slot contains the ordinary or signed local data; and
the determinant slot contains the orientation, sign, and unit conventions
needed for a single determinant line.  This decomposition is bookkeeping
only in the innocuous sense that a direct-sum decomposition is: it is
precisely the bookkeeping needed to ask whether each column is adequate.

The adequacy test is a Schur-complement calculation~\citep{Zhang2005schur,HornJohnson2013}.  Suppose that $L$ is
the native family admitted in a column and $D$ is the readout one wants to
recover from that column.  Form the block Gram quantities
\[
  \KLL,\qquad \KDL,\qquad \KDD,
\]
where $\KLL$ records the covariance of the admitted source directions,
$\KDL$ records the pairing between the readout and those source directions,
and $\KDD$ records the readout's own pre-source energy.  The adequacy
residual is
\[
  \XiC=\KDD-\KDL\,\Kdagger\,K_{LD},
\]
with $\Kdagger$ the Moore--Penrose inverse of $\KLL$ in the finite
calculation at hand.  The column is called \emph{adequate} when
$\XiC=0$: after projecting the desired readout onto the admitted source
directions, no residual direction remains.  If $\XiC>0$ in a scalar trace
calculation, then the admitted row is not enough; it has forgotten a
direction that is present in the typed target.

All adequacy computations in this paper are finite and explicit.  The
linear algebra is deliberately elementary: finite index sets, row vectors,
dot products, identity matrices, matrix multiplication, and traces.  The
collapses are not appeals to a black-box functional-analytic statement.
They are the finite identity-matrix action
\[
  I v=v
\]
and its Schur-complement consequence.  For example, if the full native
height matrix is admitted, the readout row $d_H$ for
$d(\log\det)_H$ projects through the identity source, so
$d_Hd_H^{\mathsf T}-d_H I d_H^{\mathsf T}=0$.  If only a scalar determinant
is admitted, the orthogonal complement has positive trace.  The same
pattern recurs in the analytic/period, finite-source, ordinary $p$-adic,
signed $p$-adic, and rank-two height calculations.

This finite Schur-complement language is the specialization, for the BSD
fundamental line, of the adequacy calculus developed in the surrounding
Six Birds framework \citep{TsiokosNeedles2026,TsiokosFoundationsII2026}.
For the reader of the present paper, no framework machinery is needed
beyond the following translation.  An ``adequacy residual'' is simply the
Schur complement above.  A ``typed five-column decomposition'' is simply
the decision to keep the height, period, finite, $p$-adic, and determinant
data in separate labelled columns.  An ``audit-result no-go'' is a recorded
boundary statement: it identifies where the theorem-grade literature, in
the audited form used here, does not supply the missing identity.  The
interaction vocabulary and audit discipline are inherited from the finite
calculus of \citet{TsiokosFoundationsIII2026}; the mathematical statements
below are written directly in the language of the BSD factors.

A low-conductor example makes the separation concrete.  For the curve
$11a1$~\citep{Cremona1997,LMFDB}, the scalar rank-zero BSD comparison has
\[
  \frac{L(E,1)}{\Omega_E^+}=0.2
  \qquad\text{and}\qquad
  \frac{\prod_v c_v\,\#\Sha(E/\Q)}{\#E(\Q)_{\tors}^{\,2}}
  =\frac{5}{25}=0.2.
\]
As a scalar check this is a single equality.  In the five-column atlas it is
not a single source.  The analytic/period column contributes the quotient
$L(E,1)/\Omega_E^+$; the height/regulator column is rank-zero and therefore
vacuous; the finite column remembers that the number $5/25$ arose from a
Tamagawa factor and torsion rather than from Cassels--Tate pairing data;
the local $p$-adic columns ask which ordinary or signed source data are
present at each prime; and the determinant column records the orientation
needed to assemble the product.  The scalar equality is therefore compatible
with, but coarser than, the typed comparison.  The point of the apparatus is
to measure this coarsening column by column.

This is also the reason for working with readouts rather than only with
equalities of final products.  A scalar product can be correct for several
different reasons.  In the finite factor, a valuation can sit in the
Tamagawa part or in the $\Sha$ part and still contribute the same exponent
to the public scalar quotient.  In the signed supersingular local theory,
the two signed directions can have the same unsigned aggregate while
differing in their $+$ and $-$ components.  In rank at least two, a Gram
matrix can give the regulator once a Mordell--Weil basis is already known,
but the matrix itself does not name independent motivic-cycle sources.
These are not paradoxes.  They are exactly what one should expect when a
multi-column object is projected to a single row.  The Schur residual is the
linear-algebraic device that records which directions survive the projection
and which directions have been lost.

The paper therefore uses ``native family'' in a modest sense: it means the
data one is willing to admit as the source for a column before applying the
readout.  For the analytic/period column, the native family is the pair
$(A_E,\Omega_E^+)$, where
$A_E=L^{(r_{\mathrm{an}})}(E,1)/r_{\mathrm{an}}!$ is the analytic leading
coefficient and $\Omega_E^+$ is the real period, rather than the quotient
alone.  For the finite column, it is the tower containing Tamagawa, torsion,
$\Sha$, and Cassels--Tate
pairing data rather than the scalar finite factor alone.  For the signed
$p$-adic column, it is the signed pair of local directions rather than their
unsigned sum.  The adequacy theorem for a column says that this admitted
native family is large enough for the intended readout.  A no-go theorem
says that a smaller public row is not large enough, and it identifies an
explicit pair of inputs or an explicit residual direction witnessing the
loss.

Nothing in this setup requires the reader to adopt the surrounding
terminology as a black box.  The relevant matrices are small: identity
matrices of sizes $2$, $3$, $4$, and $5$, rank-one projectors, and diagonal
traces.  The Moore--Penrose notation is used only where the finite matrix
has an evident inverse or pseudoinverse in the displayed calculation.  The
rank-one residuals are ordinary rational computations such as
\[
  \tr(I_4-P_s)=4-1=3
\]
for the scalar finite row $s=(1,-2,1,0)$, or
\[
  \tr(I_3-P_{\mathrm{uns}})=3-2=1
\]
for the unsigned signed-local projection.  Thus the reader can check the
adequacy and inadequacy claims at the same level at which the statements are
made: finite-dimensional linear algebra over explicit coordinates.

The first part of the paper constructs the five columns.  The
Bloch--Kato component map sends the integrated BSD carrier to the
determinant-line target by the direct sum of the height/regulator,
analytic/period, finite, local $p$-adic, and determinant-assembly
components.  After choosing a comparison trivialization, the corresponding
defect decomposes additively into the five component defects.  Consequently,
if every component defect vanishes in the declared comparison, then the
total Bloch--Kato defect vanishes; contrapositively, a nonzero total defect
must be witnessed by at least one nonzero or legally incomparable component.

The order of the paper follows this diagnostic structure.  The
decomposition section fixes the five columns and the additive defect
equation.  The next four sections treat the height/regulator,
analytic/period, finite-source, and $p$-adic columns separately, always
asking the same question: does the declared native family collapse the
Schur residual for the declared readout?  The determinant-assembly section
then records the one remaining orientation coordinate needed to assemble
the local tensor into a global determinant-line object.  The comparison
section contracts the typed atlas back to the usual Bloch--Kato scalar
readout and proves that the typed source is adequate for that contraction.

The middle sections prove the adequacy collapses.  The height/regulator
column admits the full Neron--Tate pairing and collapses by the identity
action.  The analytic/period column admits both logarithmic coordinates
$(\log A_E,\log\Omega_E^+)$ and collapses for the readout
$\log A_E-\log\Omega_E^+$.  The finite-source column collapses only when
the Tamagawa, torsion, $\Sha$, and Cassels--Tate coordinates are all present;
the scalar finite factor and the dimension-only $\Sha$ row leave positive
residuals.  The ordinary $p$-adic row collapses with ordinary-control data
included, while omitting that coordinate leaves a one-dimensional residual.
The signed supersingular row collapses only after separating the $+$ and
$-$ directions; unsigned aggregation forgets their difference.  The
determinant-assembly calculation then shows a structural residual exactly
in the determinant-orientation coordinate when the four constructed
components are assembled without the final determinant row.

The comparison section contracts the typed five-column tensor to the usual
Bloch--Kato scalar package.  At the log-tangent level the comparison is the
row $(1,1,1,1,0)$: it reads the height, analytic, finite, and local
$p$-adic columns and kills the determinant-orientation coordinate.  Its
Schur residual is zero against the typed source, so the typed source is
adequate for the scalar Bloch--Kato readout.  At the same time, the kernel
contains the determinant coordinate, so the scalar package is a strict
one-dimensional shadow of the typed atlas.

Two calibrations complete the apparatus.  The first is the exterior algebra
$\vsrc$, generated in rank two by symbols $\Jone,\Jtwo$ and
$\Jtwelve=\Jone\Jtwo$ with $\Jone^2=\Jtwo^2=0$ and
$\Jtwo\Jone=-\Jtwelve$.  It is a determinant-shaped source algebra: the
wedge $\Jone\wedge\Jtwo$ records independence of two virtual motivic-cycle
directions without pretending that the corresponding motivic cycles have
been constructed.  The second is the normalization constant $\kapr$.  Once
the cascade convention $c_E=\Tam(E)$ is fixed, the cascade leading-term
form is equivalent to the standard Strong-BSD leading-coefficient form
exactly when
\[
  \kapr(E)=\frac{\#\Sha(E/\Q)}{\#E(\Q)_{\tors}^{\,2}}.
\]
This is an equivalence of forms, not a derivation of BSD.

The calibration sections have a different purpose from the column
sections.  The column sections test whether a given source row is adequate
for an already specified readout.  The $\vsrc$ section asks what sort of
formal source algebra would be capable of expressing rank-two independence
before any descent to actual motivic cycles is supplied.  The
$\kapr$ section asks how the cascade normalization must be chosen so that
the cascade leading-term expression is not a new conjecture but the
standard Strong-BSD expression written in the apparatus coordinates.  These
two calibrations constrain the apparatus at opposite ends: $\vsrc$ marks
the source-side independence problem, while $\kapr$ fixes the scalar
normalization at the target side.

The final mathematical contribution is the bounded no-go suite.  The
finite scalar factor is noninjective on the typed finite tower; the
dimension of $\Sha[p]$ is a shadow of Cassels--Tate determinant data; the
Neron--Tate Gram matrix does not by itself construct rank-$r$ motivic-cycle
sources; a five-curve window cannot promote to a global theorem depending
only on those five rows; and no fixed finite prime list certifies all local
support-prime rows.  These statements are useful because they are local and
typed: each identifies a precise missing direction rather than replacing
the missing direction by an undifferentiated failure of BSD.

The scope is deliberately narrow.  The no-go suite (NG-1--NG-5,
higher-rank, global, support-prime) records where theorem-grade Mode-A
literature (the established arithmetic results audited as external
provenance) does \textbf{not} reach; it is \textbf{not} a claim that the
missing identities cannot exist.  Adequacy collapses are proved at the
finite, mathlib-free, typed level (identity-matrix action; $\Q$ rank-1
projector traces); operator/motivic full generality is not claimed.  The
$\vsrc$ exterior-product relations are represented structurally on the
$2^r$ coordinate space; the dimension/calibration content is faithful,
full axiomatization is not claimed.  The five per-column transports are
recognition sources (carriers), not derived theorems.
